\documentclass[11pt,a4paper]{article}
\usepackage[utf8]{inputenc}
\usepackage{amsmath,amssymb,amsfonts}
\usepackage{booktabs}
\usepackage{geometry}
\usepackage{hyperref}
\usepackage{microtype}

\geometry{margin=1in}

\title{\textbf{Master Equation — Liquid Tumor Model}}
\author{}
\date{}

\begin{document}

\maketitle

\section{State Space and Parameters}

The tissue has $N$ fixed sites. With a single chromosome and constant mutation rate $\mu$, the full state is the pair $(C, D)$, with $W = N - C - D$ determined implicitly.

\begin{table}[h!]
\centering
\caption{Model Parameters and Variables}
\begin{tabular}{ll}
\toprule
\textbf{Symbol} & \textbf{Meaning} \\
\midrule
$N$ & Total lattice sites (fixed) \\
$C$ & Cancer cells, division rate $r_{\max}$ \\
$W$ & Wild-type cells, division rate $r_0$ \\
$D$ & Dead cells, division rate $0$ \\
$\mu$ & Per-division mutation rate (constant) \\
$\beta$ & Moran win probability $= r_{\max} / (r_{\max} + r_0)$ \\
$P(C,D,t)$ & Probability of state $(C,D)$ at time $t$ \\
\bottomrule
\end{tabular}
\end{table}

A single chromosome means any Housekeeping (HK) mutation kills the daughter instantly (with probability $\mu$). The mother cell is \textbf{never} modified (daughter-only mutation rule).

\section{Target Selection Kernel}

For each division event of a cancer cell, a target site $j$ is drawn globally:
\begin{itemize}
    \item \textbf{With probability $D/N$}: $j$ is a dead cell $\rightarrow$ replacement always succeeds.
    \item \textbf{With probability $W/N$}: $j$ is a WT cell $\rightarrow$ Moran competition, cancer wins with probability $\beta$.
    \item \textbf{With probability $C/N$}: $j$ is a cancer cell $\rightarrow$ symmetric Moran, cancer wins with probability $1/2$.
\end{itemize}

\section{Elementary Transition Rates}

Multiplying the selection probabilities by the total cancer division rate $C \cdot r_{\max}$, we obtain the transition rates:

\begin{table}[h!]
\centering
\caption{Elementary Transitions and Rates}
\begin{tabular}{lcccl}
\toprule
\textbf{Transition} & $\Delta C$ & $\Delta D$ & \textbf{Rate } $\lambda$ \\
\midrule
Dead target, daughter lives & $+1$ & $-1$ & $C \cdot r_{\max} \cdot \frac{D}{N} \cdot (1-\mu)$ \\
Dead target, daughter dies & $0$ & $0$ & $C \cdot r_{\max} \cdot \frac{D}{N} \cdot \mu$ \quad \textit{(dead $\rightarrow$ dead, no net change)} \\
WT target, Moran win, daughter lives & $+1$ & $0$ & $C \cdot r_{\max} \cdot \frac{W}{N} \cdot \beta \cdot (1-\mu)$ \\
WT target, Moran win, daughter dies & $0$ & $+1$ & $C \cdot r_{\max} \cdot \frac{W}{N} \cdot \beta \cdot \mu$ \\
Cancer target, Moran win, daughter lives & $0$ & $0$ & $C \cdot r_{\max} \cdot \frac{C}{2N} \cdot (1-\mu)$ \quad \textit{(no net change)} \\
Cancer target, Moran win, daughter dies & $-1$ & $+1$ & $C \cdot r_{\max} \cdot \frac{C}{2N} \cdot \mu$ \\
\bottomrule
\end{tabular}
\end{table}

\begin{quote}
\textbf{Note:} When the target is dead and the daughter also dies, a dead cell is replaced by a dead cell ($\Delta D = 0$). This is a wasted division event.
\end{quote}

\section{Compact Rate Functions}

Define the four non-zero transition rates as functions of the current state $(C, D)$, with $W = N - C - D$:

\begin{align}
\lambda^+_{CD}(C,D) &= C \cdot r_{\max} \cdot \frac{D}{N} \cdot (1-\mu) && [\Delta C = +1,\ \Delta D = -1] \\
\lambda^+_{C}(C,D)  &= C \cdot r_{\max} \cdot \frac{W}{N} \cdot \beta \cdot (1-\mu) && [\Delta C = +1,\ \Delta D = 0] \\
\lambda^+_{D}(C,D)  &= C \cdot r_{\max} \cdot \frac{W}{N} \cdot \beta \cdot \mu && [\Delta C = 0,\ \Delta D = +1] \\
\lambda^-_{C}(C,D)  &= C \cdot r_{\max} \cdot \frac{C}{2N} \cdot \mu && [\Delta C = -1,\ \Delta D = +1]
\end{align}

where $\beta = \frac{r_{\max}}{r_{\max} + r_0}$ and $W = N - C - D$.

\section{Full Master Equation}

The master equation governing the probability distribution is:

\begin{align*}
\frac{\partial P(C,D,t)}{\partial t} = 
&\phantom{+} \lambda^+_{CD}(C-1,\, D+1)\cdot P(C-1,\, D+1,\, t) \tag{dead site $\to$ live cancer daughter} \\
&+ \lambda^+_{C}(C-1,\, D)\cdot P(C-1,\, D,\, t) \tag{WT site $\to$ live cancer daughter} \\
&+ \lambda^+_{D}(C,\, D-1)\cdot P(C,\, D-1,\, t) \tag{WT site $\to$ dead daughter} \\
&+ \lambda^-_{C}(C+1,\, D-1)\cdot P(C+1,\, D-1,\, t) \tag{cancer site $\to$ dead daughter} \\
&- \Big[\lambda^+_{CD}(C,D) + \lambda^+_{C}(C,D) + \lambda^+_{D}(C,D) + \lambda^-_{C}(C,D)\Big]\cdot P(C,D,t)
\end{align*}

where the last term represents the total outflow from state $(C, D)$.

\section{Mean-Field (Deterministic) Equations}

Replacing $P(C,D,t)$ with a delta function at the means $c = \langle C \rangle$, $d = \langle D \rangle$, and writing $w = N - c - d$:

\begin{equation}
\frac{dc}{dt} = c \cdot r_{\max} \left[ \frac{d}{N}(1-\mu) + \frac{w}{N}\,\beta\,(1-\mu) - \frac{c}{2N}\,\mu \right]
\end{equation}

\begin{equation}
\frac{dd}{dt} = c \cdot r_{\max} \left[ -\frac{d}{N}(1-\mu) + \frac{w}{N}\,\beta\,\mu + \frac{c}{2N}\,\mu \right]
\end{equation}

\section{Stationary Conditions}

To find stationary states, we set $\frac{dc}{dt} = 0$ and $\frac{dd}{dt} = 0$.

\subsection{Step 1: Solve $\frac{dd}{dt} = 0$ for $d^*$}

Setting the death rate change to zero:
\begin{equation}
d\,(1-\mu) = \mu\left[\beta\,(N - c - d) + \frac{c}{2}\right]
\end{equation}

Solving for $d^*$:
\begin{equation}
\boxed{d^* = \frac{\mu\left[\beta(N - c^*) + \dfrac{c^*}{2}\right]}{(1-\mu) + \mu\beta}}
\end{equation}

\subsection{Step 2: Substitute into $\frac{dc}{dt} = 0$ and solve for $c^*$}

Define two shorthand constants:
\begin{equation}
\beta \equiv \frac{r_{\max}}{r_{\max} + r_0}, \qquad \gamma \equiv \frac{(1-\mu)(1-\beta)}{(1-\mu)+\mu\beta}
\end{equation}

Substituting $d^*$ into the expression for $c^*$ yields:
\begin{equation}
\boxed{c^* = N \cdot \frac{2\beta\,[\gamma\mu + (1-\mu)]}{\mu(1-\gamma) + 2\beta\,[\gamma\mu + (1-\mu)]}}
\end{equation}

\subsection{Step 3: Stability Boundary (onset approximation, $c/N \rightarrow 0$)}

At the onset of growth, the cancer self-targeting term $c/N \rightarrow 0$ and the dead fraction is sustained by the stationary $\frac{dd}{dt} = 0$ equation alone. Setting $\frac{dc}{dt} = 0$ in the onset limit gives the \textbf{critical mutation rate}:
\begin{equation}
\boxed{\mu^*_{\text{liquid}}(p_D,\, r_{\max}) = \frac{p_D}{p_D + \beta\,(1 - p_D)}}
\end{equation}
where $p_D = D/N$ and $\beta = \frac{r_{\max}}{r_{\max}+r_0}$.

\begin{itemize}
    \item For $\mu < \mu^*$: cancer grows (gain channel dominates).
    \item For $\mu > \mu^*$: cancer shrinks (mutation-driven loss dominates).
    \item The dead-cell fraction $p_D$ \textbf{lowers} the effective threshold: a tissue with more dead cells is \textit{easier} for the tumor to invade, because dead sites are the highest-priority targets and do not fight back.
\end{itemize}

\subsection{Comparison with Solid Tumor}

\begin{table}[h!]
\centering
\caption{Comparison of Critical Mutation Rates}
\begin{tabular}{ll}
\toprule
\textbf{Model} & \textbf{Critical Mutation Rate} \\
\midrule
Solid & $\mu^*_{\text{solid}} = \frac{r_{\max}}{r_{\max} + r_0} = \beta$ \\
Liquid & $\mu^*_{\text{liquid}} = \frac{p_D}{p_D + \beta(1-p_D)}$ \\
\bottomrule
\end{tabular}
\end{table}

Since $p_D \leq 1$ and $\beta \leq 1$, one can show that $\mu^*_{\text{liquid}} \leq \mu^*_{\text{solid}}$,
i.e. the liquid tumor is \textbf{harder to eradicate} by mutational load alone.

\end{document}
